\section{导读：为什么从论文进入 AI 世界}

本节先建立这期节目的学习姿势。谢青池不是以研究员身份讲论文，而是以产品负责人、长期自学者和技术边界观察者的身份，把一年多读 200 多篇论文的路径开源出来。这个视角很有价值：它不追求把每篇论文讲成数学课，而是追问每篇论文改变了什么边界、打开了什么范式、后来被谁继承、又怎样变成产品和产业的基础。

这期的 36 篇论文可以压缩成四条主线。第一是模型范式：从 GPU、AlexNet、Attention、Transformer 到 MoE、CoT、LoRA、ReAct 和 Bitter Lesson。第二是 Infra 与数据：ZeRO（Zero Redundancy Optimizer，用优化器状态、梯度和参数分片减少训练冗余）、Scaling Law、Chinchilla、LAION、RefinedWeb、MegaScale 解释模型为什么能被放大。第三是语言模型：Word2Vec、Google Translate、GPT、BERT、GPT-3、InstructGPT、Tulu 3 解释 LLM 怎样成为今天的入口。第四是多模态：DeepVideo、双流网络、GAN、Diffusion、ViT、CLIP、Stable Diffusion、DiT 解释视觉和生成模型怎样汇入 Transformer 时代。

\begin{figure}[H]
\centering
\includegraphics[width=0.94\textwidth]{four-part-history-map.png}
\caption{36 篇论文的四条主线：模型范式、Infra/数据、语言模型和多模态共同构成 AI 变迁史。自制概念图，依据视频描述和 00:19:35--03:56:38 对谈内容整理。}
\end{figure}

\begin{knowledgebox}{读图：这不是论文清单，而是边界地图}
模型范式决定“什么结构可以学”；Infra 和数据决定“能放大到什么规模”；语言模型决定“人如何用自然语言调用智能”；多模态决定“模型如何进入视觉、视频和生成世界”。四条线互相咬合，才形成今天的 AI 产业。
\end{knowledgebox}

\begin{importantbox}{本期核心命题}
读论文的价值不是背诵论文标题，而是直接接触问题源头。真正读懂一篇关键论文，等于知道当时的研究者面对什么边界、为什么选择这个解法、后来哪一部分被时代放大、哪一部分被硬件或数据淘汰。
\end{importantbox}

\subsection{用 AI 学 AI：论文阅读闭环}

上一段解释了为什么读论文，本节看怎样把论文读下去。谢青池讲到，自己最初遇到的门槛包括数学基础、英文论文、术语多义和缺少路线图。解决办法不是硬扛，而是用 AI 做翻译、解释、追问和整理，把论文从孤立 PDF 变成可讨论、可复盘、可分享的知识节点。

\begin{figure}[H]
\centering
\includegraphics[width=0.96\textwidth]{paper-reading-loop.png}
\caption{用 AI 学 AI 的论文阅读闭环：翻译、提问、路书、复盘和分享让论文不再只是 PDF。自制概念图，依据 00:01:30--00:19:35 对谈内容整理。}
\end{figure}

\begin{knowledgebox}{读图：AI 不是替你读，而是帮你搭脚手架}
AI 可以帮助跨过英文、术语和背景知识门槛，但不能替代问题意识。好的读法是先找到问题，再让 AI 帮你翻译、解释、对比和反问，最后把论文放回历史路线图中。
\end{knowledgebox}

\subsection{本章小结}

EP117 的真正主题是“怎样进入技术世界”。36 篇论文只是载体，更深的训练是建立问题地图：知道哪些概念长期不变，哪些边界正在变化，哪些能力来自模型，哪些能力来自 Infra、数据和硬件。

